%% kapa: exam=1 round=19 year=2008 period=1 form=A subject=accounting subject_ko=회계학 seq=36 number=116
%% source: Q-Net 감정평가사 기출문제 (cst00309) artlSeq=1001962 file=extracted/1008/1001962/2008년 제19회 감정평가사 1차시험 기출문제/1교시.hwp
%% conversion: kordoc 4.24.51 -> tools/gen1.py (auto); review=proofread-sonnet; flags=none
\begin{question}{116}
(주)오성의 공정 A에서는 반제품을 생산하고 있으며, 이후 공정 B와 C로 분리되어 추가가공 된다. 공정 B를 거치면 표준제품이 생산되고, 공정 C를 거치면 고급제품이 생산된다. 20X7년 중 공정 A의 결합원가는 ￦1,200,000이었다. 공정 B의 추가가공원가는 ￦40,000이고, 공정 C의 추가가공원가는 ￦160,000이었다. 20X7년 중 표준제품은 800개가 생산되어 개당 ￦950에 판매되었고, 고급제품은 1,200개가 생산되어 개당 ￦1,200에 판매되었다면 순실현가치법에 의한 고급제품의 제조원가는 얼마인가? (단, A와 B, C공정 모두 재공품은 없다.)
\begin{choices}
  \item ￦432,000
  \item ￦778,000
  \item ￦928,000
  \item ￦945,000
  \item ￦975,000
\end{choices}
\end{question}
\answer{3}
